\paragraph{定义。}\label{sub:the-definitions}
设 $\scr{A}$ 是一个\trNonunital{}\trACategory。关于“$\scr{A}$ 具有%
\trIdentityMorphism”的含义，有若干种不同的版本。
\begin{itemize}
\item 称 $\scr{A}$ 是\textbf{\trStrictlyUnital}（strictly unital）的，若对每个%
\trObject{} $X$ 都存在一个（必然是唯一的）$e_X\in\hom^0_{\scr{A}}(X,X)$，满足
\begin{equation}\label{eq:strictly-unital}
\begin{aligned}
&\mu^1_{\scr{A}}(e_X)=0;\\
&(-1)^{|a|}\mu^2_{\scr{A}}(e_{X_1},a)=a=\mu^2_{\scr{A}}(a,e_{X_0})
\quad\text{其中 }a\in\hom_{\scr{A}}(X_0,X_1);\\
&\mu^d_{\scr{A}}(a_{d-1},\ldots,a_{n+1},e_{X_n},a_n,\ldots,a_1)=0\\
&\text{其中 }d>2,\ a_k\in\hom_{\scr{A}}(X_{k-1},X_k)\text{，而 }0\le n<d
\text{ 任意}.
\end{aligned}
\end{equation}
\item 称 $\scr{A}$ 是\textbf{\trCunital}（cohomologically unital or c-unital）的，%
若 $H(\scr{A})$ 是\trUnital{}的（即对任何\trObject{} $X$ 都有\trIdentityMorphism
——这正意味着它是一个通常意义下的\trGraded\trLinear\trCategory）。
\item 一个\textbf{\trHomotopy\trUnital}（homotopy unital）的\trACategory{}是指%
集合 $\Ob\scr{A}$，连同诸\trGraded\trVectorSpace{} $\hom_{\scr{A}}(X_0,X_1)$ 与%
诸\trMultilinear{}映射
\begin{equation}\label{eq:homotopy-unital-maps}
\begin{aligned}
&\mu^{d,(i_d,\ldots,i_0)}_{\scr{A}}\colon
\hom_{\scr{A}}(X_{d-1},X_d)\otimes\cdots\otimes\hom_{\scr{A}}(X_0,X_1)\\
&\hspace{6em}\longrightarrow
\hom_{\scr{A}}(X_0,X_d)[2-d-2{\textstyle\sum_k} i_k]
\end{aligned}
\end{equation}
（对一切 $d+i_0+\cdots+i_d>0$）。这些映射必须满足某些广义的结合性方程；当
$i_0=\cdots=i_d=0$ 时，它们\trDegenerate{}为\eqref{eq:aoo-associativity}。最先出现的几条
新方程是
\begin{equation}\label{eq:homotopy-unital-first-equations}
\begin{aligned}
&\mu^1_{\scr{A}}(\mu^{0,(1)}_{\scr{A}})=0,\\
&(-1)^{|a|-1}\mu^2_{\scr{A}}(\mu^{0,(1)}_{\scr{A}},a)
+\mu^1_{\scr{A}}(\mu^{1,(1,0)}_{\scr{A}}(a))
+\mu^{1,(1,0)}_{\scr{A}}(\mu^1_{\scr{A}}(a))=a,\\
&\mu^2_{\scr{A}}(a,\mu^{0,(1)}_{\scr{A}})
+\mu^1_{\scr{A}}(\mu^{1,(0,1)}_{\scr{A}}(a))
+(-1)^{|a|-1}\mu^{1,(0,1)}_{\scr{A}}(\mu^1_{\scr{A}}(a))=-a,\\
&\mu^{1,(1,0)}_{\scr{A}}(\mu^{0,(1)}_{\scr{A}})
+\mu^{1,(0,1)}_{\scr{A}}(\mu^{0,(1)}_{\scr{A}})
+\mu^1_{\scr{A}}(\mu^{0,(2)}_{\scr{A}})=0.
\end{aligned}
\end{equation}
撇开通常的正负号不谈，此式的前两部分说的是：用\trCocycle
$e_X=-\mu^{0,(1)}_{\scr{A}}$ 在任一侧作乘法，都\trChainHomotopy{}于恒等。利用
这两个\trChainHomotopy{}中的任何一个，可以看出 $\mu^2_{\scr{A}}(e_X,e_X)$ 等于
$e_X$ 加一个\trCoboundary。这两个\trCoboundary{}之差是一个\trCocycle，但
\eqref{eq:homotopy-unital-first-equations} 的最后一部分表明，它在\trCohomology{}%
中本身就是平凡的。整组方程可以想成是把这个层层升高的\trHomotopy{}等级结构不断
延续下去；细节见 [55, Section 5]。
\end{itemize}

任何\trStrictlyUnital{}的\trACategory{}都是\trHomotopy\trUnital{}的（取
$\mu^{0,(1)}_{\scr{A}}=-e_X$，并把其余满足 $i_0+\cdots+i_d>0$ 的映射
\eqref{eq:homotopy-unital-maps} 全取为零）。此外，正如上面的讨论所表明的，
\trHomotopy\trUnital{}性蕴含\trCunital{}性。

我们还应当提到\trFunctor{}方面的相应概念。设 $\scr{A}$、$\scr{B}$ 是\trStrictlyUnital{}的，则\trStrictlyUnital{}的%
\trAFunctor{} $\scr{F}\colon\scr{A}\to\scr{B}$ 是指满足
$\scr{F}^1(e_X)=e_{\scr{F}(X)}$（对一切\trObject{} $X$）以及
$\scr{F}^d(a_{d-1},\ldots,a_{n+1},e_{X_n},a_n,\ldots,a_1)=0$（对任意 $d\ge 2$ 与%
任意 $n$）的\trAFunctor。若 $\scr{A}$、$\scr{B}$ 是\trCunital{}的，则 \trCunital{}%
的\trAFunctor{} $\scr{F}\colon\scr{A}\to\scr{B}$ 简单地说就是使 $H(\scr{F})$ 为%
\trUnital{}的那一种。\trHomotopy\trUnital{}版本见 [55]。

\trStrictlyUnital{}性在许多方面都很方便，也因此被文献广泛采用，尽管它看起来似乎与\trHomotopy{}代数
结构的一般哲学不合拍。\trCunital{}性\trApriori{}看也可能显得可疑，因为它落到了
另一个极端，把这个问题整个下放到\trCohomology{}层面。就我们预期的应用而言，
\trFukayaCategory{}并不天然带有\trStrictlyUnital{}元，但容易看出它们是\trCunital{}的。事实上，
对它们存在一种\trHomotopy\trUnital{}元的几何构造，它用补充的\trModuliSpace{}来
定义所需的满足 $i_0+\cdots+i_d>0$ 的映射\eqref{eq:homotopy-unital-maps}，见
[60, \S20]。虽然这一切似乎都指向把\trHomotopy\trUnital{}性当作最合适的概念，
接下来的讨论却要出现一个可能出人意料的转折：对相关\trObstruction{}的考察表明，
存在一个一般的代数程序，能为任何\trCunital{}的\trACategory{}配备%
\trHomotopy\trUnital{}元。此外，从 [55] 中的推导可以看出，任何%
\trHomotopy\trUnital{}的\trACategory{}都与某个\trStrictlyUnital{}的\trACategory{}
\trQuasiisomorphism。这意味着这三种概念之间的差别属于口味与便利的问题，而非
实质问题。我们的选择是主要在\trCunital{}的世界里工作。不过，作为证明某些基本
性质时的技术工具，我们将给出一种把\trCunital{}的\trACategory{}变成\trStrictlyUnital{}%
\trACategory{}的直接方法。

\begin{lemma}\label{thm:strictly-unital-modification}
设 $\scr{A}$ 是\trCunital{}的\trACategory。那么存在一个\trFormalDiffeomorphism{}%
$\Phi$，$\Phi^1=\mathrm{Id}$，使得改动后的 \trAinfty-结构 $\Phi_*\scr{A}$ 是%
\trStrictlyUnital{}的。
\end{lemma}

\begin{proof}
最初的一步如下。对每个\trObject{} $X$，选取一个在 $H(\scr{A})$ 中代表恒等的%
\trCocycle{} $e_X\in\hom^0_{\scr{A}}(X,X)$。对任意三个\trObject，都能找到一个%
\trChainMap{}%
$\mu^2_{\tilde{\scr{A}}}\colon\hom_{\scr{A}}(X_1,X_2)\otimes\hom_{\scr{A}}(X_0,X_1)
\to\hom_{\scr{A}}(X_0,X_2)$，它代表 $H(\scr{A})$ 中的\trComposition，且满足
\eqref{eq:strictly-unital} 中的第二个条件。取一个满足 $\Phi^1=\mathrm{Id}$ 的%
\trFormalDiffeomorphism，使得每个 $\Phi^2$ 都是介于 $\mu^2_{\scr{A}}$ 与
$\mu^2_{\tilde{\scr{A}}}$ 之间的一个\trChainHomotopy。于是
$\tilde{\scr{A}}=\Phi_*\scr{A}$ 的前两个\trComposition{}映射即为
$\mu^1_{\tilde{\scr{A}}}=\mu^1_{\scr{A}}$ 与所选的 $\mu^2_{\tilde{\scr{A}}}$。

从这里开始，构造成为递归的。设有\trNonunital{}\trACategory{} $\scr{A}$，连同%
满足\eqref{eq:strictly-unital} 中前两个方程的元素 $e_X\in\hom^0_{\scr{A}}(X,X)$。
考虑对这种 $\scr{A}$ 的下面这族条件 $(U_{d,n})$，其标号由 $d>2$ 与 $0\le n\le d$
给出：
\begin{equation*}
\begin{aligned}
&(U_{d,n})\ \ \mu^i_{\scr{A}}(a_{i-1},\ldots,a_{j+1},e_{X_j},a_j,\ldots,a_1)
\text{\kaishu{} 在下列两种情形下消没：}\\
&\hspace{6em}\text{\kaishu{}对 }i<d\text{\kaishu{} 而 }j\text{\kaishu{} 任意，%
以及 }i=d\text{\kaishu{} 而 }j<n\text{\kaishu{}。}
\end{aligned}
\end{equation*}
$n=0$ 的情形纳入其中纯粹是出于记号上的原因，因为显然
$(U_{d,0})=(U_{d-1,d-1})$。设 $\scr{A}$ 已对某个 $n<d$ 满足 $(U_{d,n})$。取一个%
\trFormalDiffeomorphism{} $\Phi$，使得 $\Phi^1=\mathrm{Id}$，$\Phi^k=0$（对
$2\le k<d-1$），并且
\begin{equation}\label{eq:formal-diffeomorphism-recursion}
\begin{gathered}
\Phi^{d-1}(a_{d-1},\ldots,a_1)=(-1)^{\clubsuit_n}\,\mu^d_{\scr{A}}
(a_{d-1},\ldots,a_{n+1},e_{X_n},a_n,\ldots,a_1),\\
\Phi^d(a_d,\ldots,a_1)=(-1)^{\clubsuit_n}\,\mu^{d+1}_{\scr{A}}
(a_d,\ldots,a_{n+1},e_{X_n},a_n,\ldots,a_1).
\end{gathered}
\end{equation}
由定义显然可见，$\tilde{\scr{A}}=\Phi_*\scr{A}$ 中的\trComposition{}映射在阶
$<d-1$ 时与 $\scr{A}$ 的一致。稍复杂一点的一个计算（用到 \trAinfty-结合性方程
与归纳假设）表明，同一结论在阶 $d-1$ 处也成立。现在看阶 $d$。为简化公式，记
$(b_1,\ldots,b_{d+1})=(a_1,\ldots,a_n,e_{X_n},a_{n+1},\ldots,a_d)$。可以算得
\begin{equation}\label{eq:modified-composition}
\begin{aligned}
&\mu^d_{\tilde{\scr{A}}}(a_d,\ldots,a_1)\\
&\hspace{1.5em}=\mu^d_{\scr{A}}(a_d,\ldots,a_1)
-(-1)^{\text{\S}}\,\mu^2_{\scr{A}}\bigl(\mu^d_{\scr{A}}
(a_d,\ldots,e_{X_{n+1}},a_{n+1},\ldots,a_2),a_1\bigr)\\
&\hspace{3em}-(-1)^{\clubsuit_{n+1}}\sum_{j<n}(-1)^{\dagger}\,\mu^d_{\scr{A}}
\bigl(a_d,\ldots,e_{X_{n+1}},a_{n+1},\ldots\\
&\hspace{14em}\ldots,\mu^2_{\scr{A}}(a_{j+2},a_{j+1}),\ldots,a_1\bigr)\\
&\hspace{3em}-(-1)^{\clubsuit_n}\sum_{i,j}(-1)^{\ddagger}\,
\mu^{d+2-i}_{\scr{A}}
\bigl(b_{d+1},\ldots,\mu^i_{\scr{A}}(b_{j+i},\ldots,b_{j+1}),b_j,\ldots,b_1\bigr)
\end{aligned}
\end{equation}
其中正负号为 $\text{\S}=|a_2|+\cdots+|a_{n+1}|-n$，
$\dagger=|a_1|+\cdots+|a_j|-j$，$\ddagger=|b_1|+\cdots+|b_j|-j$，而最后一个求和
取遍下面四种可能：$i=1$ 且 $j\ne n$；$i=2$ 且 $j>n$；$i=d$ 且 $j=0$；最后是
$i=d+1$。注意，最后提到的这个求和中的表达式，在 $i=1$、$j=n$ 时，或当
$3\le i\le d-1$ 时，也由归纳假设而消没。

先设 $n=0$，此时\eqref{eq:modified-composition} 右边的第三个表达式是一个空求和。
把它去掉，再对第四个表达式应用 \trAinfty-结合性方程，便得到
\begin{equation*}
\mu^d_{\tilde{\scr{A}}}(a_d,\ldots,a_1)=\mu^d_{\scr{A}}(a_d,\ldots,a_1)
-\mu^2_{\scr{A}}\bigl(\mu^d_{\scr{A}}(a_d,\ldots,a_2,e_{X_1}),a_1\bigr).
\end{equation*}
若 $a_1$ 也是恒等元，此式显然消没，这说明 $\tilde{\scr{A}}$ 具有性质
$(U_{d,1})$。接着考虑 $n>0$ 的情形。此时，若取 $i=2$ 并让求和取遍一切 $j\le n$，
\eqref{eq:modified-composition} 中最后一个表达式里的被加项也消没；对 $i=d$ 而
$j=1$ 亦然。与前一样地运用 \trAinfty-结合性方程，即可证明这个求和消没。把这一项
消灭之后，容易验证 $\tilde{\scr{A}}$ 具有性质 $(U_{d,n+1})$。

反复进行这一过程，就得到一列\trFormalDiffeomorphism{} $\Phi_1,\Phi_2,\ldots$，
使得当 $k\to\infty$ 时，$(\Phi_k\circ\cdots\circ\Phi_1)_*\scr{A}$ 满足%
\trStrictlyUnital{}条件中越来越多的部分。诸 $\Phi_k$ 本身则与恒等越来越贴近，
意思是说，随着 $k\to\infty$，$\Phi_k$ 的第一个非平凡项将出现在越来越高的阶上。
这蕴含无穷\trComposition{} $\Phi=\cdots\circ\Phi_3\circ\Phi_2\circ\Phi_1$ 有意义，并且
$\Phi_*\scr{A}$ 是\trStrictlyUnital{}的，正如所要证。
\end{proof}

\begin{remark}\label{thm:strictly-unital-history}
\cref{thm:strictly-unital-modification} 由 Lefevre [95, Theorem 3.2.1.1] 与本书
作者各自独立得到。存在一个更概念性的证明，它基于\trHomology{}\trPerturbation{}理论，以及
“幺代数的完整\trHochschild\trComplex{}与既约\trHochschild\trComplex{}互为%
\trQuasiisomorphism”这一事实（\eqref{eq:formal-diffeomorphism-recursion} 中的
诸表达式，正是受这些\trComplex{}之间的经典\trHomotopyEquivalence{}启发而来）。类似的
结论对 \trAFunctor{} [95, Theorem 3.2.2.1]、\trAinfty-\trModule{} [95,
Theorem 3.3.1.2] 等等也成立。
\end{remark}
